\documentclass[10pt,a4paper]{amsart}
\usepackage[margin=19mm]{geometry}
\usepackage{amsmath,amssymb,mathtools}
\usepackage[scr=rsfs,cal=euler]{mathalfa}
\usepackage{xfrac}
\usepackage[unicode,hidelinks,bookmarks=false]{hyperref}
\newcommand{\ie}{i.e.\ }
\newcommand{\cf}{cf.\ }
\newcommand{\resp}{respectively\ }
\newcommand{\Subsection}{\subsection}
\providecommand{\nobreakdash}{\nobreak\hbox{-}}
\setlength{\emergencystretch}{2em}
\multlinegap0pt
\usepackage{bm}
\usepackage{bbm}
\usepackage{dsfont}
\usepackage{xspace}
\usepackage{cancel}
\usepackage{mathtools}
\usepackage{amsthm}
\makeatletter
\@ifundefined{theorem}{\newtheorem{theorem}{Theorem}}{}
\@ifundefined{lemma}{\newtheorem{lemma}[theorem]{Lemma}}{}
\makeatother
\def\R{\mathbb{R}}
\newcommand{\bu}{\mathbf u}
\newcommand{\bv}{\mathbf v}
\newcommand{\bw}{\mathbf w}
\newcommand{\bx}{\mathbf x}
\newcommand{\by}{\mathbf y}
\newcommand{\ones}{{\bf{1}}}
\newcommand{\zero}{{\bf{0}}}
\title[Perfect state transfer between real pure states]{Perfect state transfer between real pure states}
\author{Chris Godsil, Stephen Kirkland, Hermie Monterde}
\date{Source: arXiv:2502.08103v4 (CC BY 4.0)}
\begin{document}
\maketitle

\noindent\emph{Source and licence.} Perfect state transfer between real pure states. Version arXiv:2502.08103v4 (CC BY 4.0), distributed under the Creative Commons Attribution 4.0 International licence, as recorded in the arXiv licence metadata for this exact version and shown on the abstract page https://arxiv.org/abs/2502.08103v4. Full rights notice: licenses/arxiv-2502.08103-CC-BY-4.0.txt.\par\smallskip

\noindent\textit{Editorial scope.} Counted unit: the Theorem whose source label is \texttt{Thm:prescribedPST} in the article (source0.tex lines 526--529, source label \texttt{Thm:prescribedPST}), delivered together with the article's own complete original proof of it (source0.tex lines 531--563, one \texttt{proof} environment of 33 lines). No mathematics is written, repaired or re-derived: statement and proof are the article's own text line for line. Required context retained uncounted: Lem:minperiod (lines 276--282); Prop:sc (lines 370--383); Thm:PSTchar (lines 434--446). Every definition, display and label of those lines is retained unchanged. Omitted from this package and not redistributed: 1--275, 283--369, 384--433, 447--525, 530--530, 564--1130. Nothing inside a retained excerpt is omitted or altered: every retained body is delivered exactly as the article's lines are written, so no omission receipt of any kind accompanies this package. Citations: the retained excerpts carry no citation command at all, so no bibliography entry of the article is transcribed and no reference is redistributed. Reference adaptation: none was needed and none was made; every reference inside the delivered unit resolves inside it, so no unresolved reference or citation is produced. Printed slips kept literal: no printed word, symbol, hypothesis or number of the counted statement or of its proof was altered. Layout and class: the article's own class is \texttt{article} with packages including mathpazo, amsfonts, amsmath, graphicx, latexsym, mathabx, geometry, fontenc, times, color, array, enumerate, amssymb, amsthm; the wrapper is the lane's pinned \texttt{amsart} editorial wrapper at 10pt on A4, which restates the theorem-like environments the retained text needs and adds the article's own macro definitions that the retained text needs (\texttt{\textbackslash R}, \texttt{\textbackslash bu}, \texttt{\textbackslash bv}, \texttt{\textbackslash bw}, \texttt{\textbackslash bx}, \texttt{\textbackslash by}, \texttt{\textbackslash ones}, \texttt{\textbackslash zero}), with extra wrapper packages (bm, bbm, dsfont, xspace, cancel, mathtools). The delivered numbering is the unit's own. Assets: no retained span contains a figure environment, a graphics-inclusion command, a PDF-page inclusion, a table or a TikZ picture; no image file is copied into this package, the package declares no asset dependency and no image file is redistributed with it. Licence: version arXiv:2502.08103v4 of this article is distributed under the Creative Commons Attribution 4.0 International licence, as recorded in the arXiv licence metadata for this exact version and shown on the abstract page https://arxiv.org/abs/2502.08103v4. A full rights notice is delivered at licenses/arxiv-2502.08103-CC-BY-4.0.txt. Characters: every retained excerpt is pure ASCII, so the delivered engine, XeLaTeX with its default Latin Modern text font, reports no missing character.
\begin{lemma}\label{Lem:minperiod}
Let $G$ be a graph and $\bx\in\R^n\backslash\{\zero\}$ with $\sigma_{\bx}(M)=\{\lambda_1,\ldots,\lambda_m\}$, where $\lambda_1>\lambda_2$.
\begin{enumerate}
\item If $m=2$, then $\bx$ is periodic in $G$ with $\rho=\frac{2\pi}{\lambda_1-\lambda_2}$
\item If $m\geq 3$ and $\bx$ is periodic in $G$, then $\rho=\frac{2\pi q}{\lambda_1-\lambda_2}$, where $q=\operatorname{lcm}(q_3,\ldots,q_{m})$ and the $p_j$'s and $q_j$'s are integers such that $\frac{\lambda_1-\lambda_j}{\lambda_1-\lambda_2}=\frac{p_j}{q_j}$ and $\operatorname{gcd}(p_j,q_j)=1$.
\end{enumerate}
\end{lemma}

\begin{theorem}
\label{Prop:sc}
Let $\bx,\by\in \R^n\backslash\{\zero\}$ and suppose
$\bx=\sum_{j\in \sigma_{\bx}(M)}\bu_j$, where each $\bu_j$ is a real eigenvector associated with an eigenvalue $\lambda_j\in \sigma_{\bx}(M)$. The following are equivalent.
\begin{enumerate}
\item The vectors $\bx$ and $\by$ are strongly cospectral.
\item There exists an orthogonal matrix $Q$ that is a polynomial in $M$ such that $Q^2=I$ and $Q\bx=\by$.
\item For some nonempty sets $\sigma_1$ and $\sigma_2$ that partition $\sigma_{\bx}(M)$, we have
\begin{equation*}
\bx=\sum_{\lambda_j\in \sigma_1}\bu_j+\sum_{\lambda_j\in \sigma_2}\bu_j\quad \text{and}\quad \by=\sum_{\lambda_j\in \sigma_1}\bu_j-\sum_{\lambda_j\in \sigma_2}\bu_j.
\end{equation*}
In this case, $\sigma_{\bx,\by}^{+}(M)=\sigma_1$ and $\sigma_{\bx,\by}^{-}(M)=\sigma_2$.
\end{enumerate}
\end{theorem}

\begin{theorem}
\label{Thm:PSTchar}
Let $G$ be a graph on $n$ vertices and let $\bx,\by \in \R^n\backslash\{\zero\}$ with $\by \ne \pm \bx$.
\begin{enumerate}
\item If $|\sigma_{\bx}(M)|=2$, then $\bx$ and $\by$ admit perfect state transfer if and only if they are strongly cospectral.
\item Let $\sigma_{\bx}(M)=\{\lambda_1,\ldots,\lambda_m\}$ for some $m\geq 3$. The vectors $\bx$ and $\by$ admit perfect state transfer if and only if $\bx$ and $\by$ are strongly cospectral, $\frac{\lambda_1-\lambda_j}{\lambda_1-\lambda_2}=\frac{p_j}{q_j}$ for each $j\geq 3$, where $p_j$ and $q_j$ are integers such that $\operatorname{gcd}(p_j,q_j)=1$, and one of the following conditions holds.
\begin{enumerate}
\item If $\lambda_1,\lambda_2\in\sigma_{\bx,\by}^+(M)$, then $$\nu_2(q_{\ell})=\nu_2(q_{k})>\nu_2(q_h)$$ for all $\lambda_{\ell},\lambda_k\in \sigma_{\bx,\by}^-(M)$ and $\lambda_h\in \sigma_{\bx,\by}^+(M)\backslash\{\lambda_1,\lambda_2\}$, where each $\nu_2(q_h)$ above is absent whenever $|\sigma_{\bx,\by}^+(M)|=2$.
\item If $\lambda_1\in\sigma_{\bx,\by}^+(M)$ and $\lambda_2\in\sigma_{\bx,\by}^-(M)$, then each $q_j$ is odd, and $p_h$ is even if and only if $\lambda_h\in\sigma_{\bx,\by}^+(M)$.
\end{enumerate}
\end{enumerate}
Moreover, the minimum PST time is $\frac{\rho}{2}$, where $\rho$ is given in Lemma \ref{Lem:minperiod}.
\end{theorem}

\begin{theorem}
\label{Thm:prescribedPST}
Let $\bx,\by\in\R^n\backslash\{\zero\}$ with $\by\neq\pm \bx$. For all $\tau>0$ and for all integers $m_1,m_2\geq 1$ such that $m_1+m_2\leq n$, there exists a real symmetric matrix $M$ such that perfect state transfer occurs between $\bx$ and $\by$ relative to $M$ at time $\tau$, $|\sigma_{\bx,\by}^+(M)|=m_1$ and $|\sigma_{\bx,\by}^-(M)|=m_2$.
\end{theorem}

\begin{proof}
Let $\bx,\by\in\R^n$ with $\by\neq\pm \bx$. Fix $\tau>0$ and fix integers $m_1,m_2\geq 1$ such that $m_1+m_2\leq n$. Since $\|\bx\|=\|\by\|$, it follows that $\bx+\by$ and $\bx-\by$ are orthogonal vectors. Since
\begin{center}
$\{\bu_1,\bu_2\}:=\left\{\frac{1}{\sqrt{m_1}}\left[\begin{array}{ccc} \ones_{m_1} \\ \zero_{m_2} \\ \zero_{n-(m_1+m_2)} \end{array} \right],\frac{1}{\sqrt{m_2}}\left[\begin{array}{ccc} \zero_{m_1} \\ \ones_{m_2} \\ \zero_{n-(m_1+m_2)} \end{array} \right]\right\}\quad $ and $\quad \{\bv_1,\bv_2\}:=\left\{\frac{\bx+\by}{\|\bx+\by\|},\frac{\bx-\by}{\|\bx-\by\|}\right\}$    
\end{center}
are orthonormal sets, we may extend them to orthonormal bases $\mathcal{U}=\{\bu_1,\ldots,\bu_n\}$ and $\mathcal{V}=\{\bv_1,\ldots,\bv_n\}$ for $\R^n$, respectively. Then there exists an orthogonal matrix $Q$ such that $Q\bu_j=\bv_j$ for each $j$. If we write $Q=[\bw_1,\ldots,\bw_n]$, where $\bw_1,\ldots,\bw_n$ are columns of $Q$, then $Q\bu_1=\bv_1$ and $Q\bu_2=\bv_2$ are equivalent to
\begin{equation*}
\bx+\by=\displaystyle\frac{\|\bx+\by\|}{\sqrt{m_1}}\left(\sum_{j=1}^{m_1}\bw_j\right)\quad \text{and} \quad \bx-\by=\displaystyle\frac{\|\bx-\by\|}{\sqrt{m_2}}\left(\sum_{j=m_1+1}^{m_2}\bw_j\right).
\end{equation*}
Thus, we obtain
\begin{equation}
\label{Eq:sc}
\bx=\displaystyle\left(\sum_{j=1}^{m_1}\displaystyle\frac{\|\bx+\by\|}{2\sqrt{m_1}}\bw_j\right)+\left(\sum_{j=m_1+1}^{m_2}\displaystyle\frac{\|\bx-\by\|}{2\sqrt{m_2}}\bw_j\right)\quad \text{and} \quad \by=\displaystyle\left(\sum_{j=1}^{m_1}\displaystyle\frac{\|\bx+\by\|}{2\sqrt{m_1}}\bw_j\right)-\left(\sum_{j=m_1+1}^{m_2}\displaystyle\frac{\|\bx-\by\|}{2\sqrt{m_2}}\bw_j\right)
\end{equation}
Now, let 
$\{\theta_1,\ldots,\theta_{n}\}$ be a set of $n$ distinct real numbers and consider the real symmetric matrix
\begin{equation*}
M=\sum_{j=1}^{n}\theta_j\bw_j\bw_j^\top
\end{equation*}
Since equation (\ref{Eq:sc}) holds, Theorem \ref{Prop:sc}(3) implies that $\bx$ and $\by$ are strongly cospectral relative to the matrix $M$ with $\sigma_{\bx,\by}^+(M)=\{\theta_1,\ldots,\theta_{m_1}\}$ and $\sigma_{\bx,\by}^-(M)=\{\theta_{m_1+1},\ldots,\theta_{m_1+m_2}\}$. We proceed with cases.

\noindent \textbf{Case 1.} Suppose $m_1+m_2=2$. In this case, choose $\theta_1,\theta_2$ such that $\theta_1-\theta_2=\frac{\pi}{\tau}$. Since $|\sigma_{\bx}(M)|=2$, invoking Theorem \ref{Thm:PSTchar}(1) yields PST between $\bx$ and $\by$ relative to $M$ at time $\frac{\pi}{\theta_1-\theta_2}=\frac{\pi}{\pi/\tau}=\tau$.

\noindent \textbf{Case 2.} Suppose $m_1+m_2\geq 3$. In this case, choose $\theta_1,\ldots,\theta_{m_1+m_2}$ such that $\theta_1-\theta_j=\pi b_j/(g\tau)$ where $b_j$ is even for all $j\in\{2,\ldots,m_1\}$, $b_j$ is odd otherwise, and $g=\operatorname{gcd}(b_2,\ldots,b_{m_1+m_2})$. Observe that
\begin{equation}
\label{Eq:PST}
\frac{\theta_1-\theta_j}{\theta_1-\theta_2}=\frac{b_j/g_j}{b_2/g_j}:=\frac{p_j}{q_j},
\end{equation}
where each $g_j=\operatorname{gcd}(b_j,b_2)$.
Thus, $q=\operatorname{lcm}(q_3,\ldots,q_{m_1+m_2})=\operatorname{lcm}(\frac{b_2}{g_3},\ldots,\frac{b_2}{g_{m_1+m_2}})=\frac{b_2}{g}$. We have two cases. First, if $m_1=1$, then each $p_j$ and $q_j$ is odd. In this case, invoking Theorem \ref{Thm:PSTchar}(2b) yields PST between $\bx$ and $\by$ relative to $M$ at time $\frac{\pi q}{\theta_1-\theta_2}=\frac{\pi b_2/g}{\pi b_2/ (g\tau)}=\tau$. Now, if $m_1\geq 2$, then $b_{\ell}$ is odd for each $\theta_{\ell}\in\sigma_{\bx,\by}^-(M)$, and so $\nu_2(g_{\ell})=\nu_2(\operatorname{gcd}(b_{\ell},b_2))=0$. Consequently,  $\nu_2(q_{\ell})=\nu_2(\frac{b_2}{g_{\ell}})=\nu_2(b_2)$ for each $\theta_{\ell}\in\sigma_{\bx,\by}^-(M)$. Moreover, $b_{h}$ is even for each $\theta_{h}\in\sigma_{\bx,\by}^+(M)$, and so $\nu_2(g_{h})=\nu_2(\operatorname{gcd}(b_{h},b_2))\geq 1$. Consequently,  $\nu_2(q_{h})=\nu_2(\frac{b_2}{g_{h}})<\nu_2(b_2)$ for each $\theta_{h}\in\sigma_{\bx,\by}^+(M)$. Equivalently, $$\nu_2(q_{\ell})=\nu_2(q_{k})>\nu_2(q_h)$$ for all $\theta_{\ell},\theta_k\in \sigma_{\bx,\by}^-(M)$ and $\theta_h\in \sigma_{\bx,\by}^+(M)\backslash\{\theta_1,\theta_2\}$, where each $\nu_2(q_h)$ above is absent whenever $|\sigma_{\bx,\by}^+(M)|=2$. Finally, invoking Theorem \ref{Thm:PSTchar}(2a) yields PST between $\bx$ and $\by$ at time $\tau$.

Combining the above two cases yields the desired result.
\end{proof}

\end{document}
